% 原书物理页=印刷页20；2.4 全节.
\section{若干经典群}\label{sec:fg-2-4}
除了第 I 章中已经遇到的例子，如循环群 $C_n=\mathbb{Z}/n\mathbb{Z}$、对称群与交错群 $S_n,A_n$，以及矩阵群 $\mathrm{GL}_n(k)$ 的一些子商之外，还可以举出：
\begin{numpara}[模群]\label{np:fg-2-4-1}
它是 $\mathrm{PSL}_2(\mathbb{R})$ 中阶为无穷的子群 $\mathrm{PSL}_2(\mathbb{Z})=\mathrm{SL}_2(\mathbb{Z})/\{\pm 1_2\}$. 群 $\mathrm{SL}_2(\mathbb{Z})$ 由矩阵
\[
S=\begin{pmatrix}0 & -1\\ 1 & 0\end{pmatrix}
\]
（阶为 $4$）与
\[
U=\begin{pmatrix}0 & -1\\ 1 & 1\end{pmatrix}
\]
（阶为 $6$）生成. 群 $\mathrm{PSL}_2(\mathbb{Z})$ 以 $(\{a,b\}\mid\{a^2,b^3\})$ 为表述.
\end{numpara}
\begin{numpara}[$n$-传递群]\label{np:fg-2-4-2}
设 $E$ 是一个集合，$n$ 是满足 $n\leqslant\operatorname{card}(E)$ 的整数，$G$ 是对称群 $S_E$ 的一个子群. 称 $G$ 是 $n$-\newterm{fg-n-transitive}{传递的}{transitive}（或者更恰当地说：$G$ 在 $E$ 上 $n$ 重传递地作用），如果对 $E$ 中任意两个由 $n$ 个元素组成的子集 $\{x_1,\ldots,x_n\}$、$\{y_1,\ldots,y_n\}$，都存在元素 $g\in G$，使得 $g.x_i=y_i$ 对一切 $i=1,\ldots,n$ 成立. 按定义，$G$ 是 $1$-传递的，就是说 $G$ 在 $E$ 上的作用是传递的. 群 $S_n$ 是 $n$-传递的，因此 $A_n$ 是 $(n-2)$-传递的（若 $n\geqslant 3$）. 对任何域 $k$，群 $\mathrm{PGL}_2(k)$（相应地 $\mathrm{PSL}_2(k)$）在 $\mathbb{P}_1(k)$ 上 $3$ 重（相应地 $2$ 重）传递地作用. 而当 $k\neq\mathbb{F}_3$ 时，$\mathrm{PGL}_2(k)$ 不是 $4$-传递的：这是因为四元组的\newterm{fg-cross-ratio}{交比}{cross-ratio} $(x_1:x_2:x_3:x_4):=\frac{x_3-x_1}{x_3-x_2}:\frac{x_4-x_1}{x_4-x_2}$ 在它的作用下不变，并且 $(\infty:0:1:x)=x$. 除 $A_n$、$S_n$ 本身之外，人们所知的 $4$ 或 $5$-传递群很少（Mathieu 群），并且猜测没有 $6$-传递的群.
\end{numpara}
下面来看几个新的矩阵型群.
\begin{numpara}[辛群与正交群]\label{np:fg-2-4-3}
设 $k$ 是交换域，$n>0$ 是整数，$J=\begin{pmatrix}0_n & 1_n\\ -1_n & 0_n\end{pmatrix}$ 是表示向量空间 $k^{2n}$ 上的交错双线性型 $(\underline{x},\underline{y})\mapsto\sum_{i=1,\ldots,n}x_iy_{n+i}-x_{n+i}y_i$ 的 $2n$ 阶方阵.\newterm{fg-symplectic-group}{辛群}{symplectic group}是群 $\mathrm{Sp}_{2n}(k)=\{P\in\mathrm{GL}_{2n}(k),\ {}^tPJP=J\}$. 于是 $\mathrm{Sp}_2(k)=\mathrm{SL}_2(k)$. 更一般地，可以证明辛矩阵的行列式总等于 $1$，从而对一切 $n$，$\mathrm{Sp}_{2n}(k)$ 是 $\mathrm{SL}_{2n}(k)$ 的子群.

\newterm{fg-orthogonal-group}{正交群}{orthogonal group}是群 $\mathrm{O}_n(k)=\{P\in\mathrm{GL}_n(k),\ {}^tP.P=\mathbf{1}_n\}$. 它的元素表示空间 $k^n$ 的``正交变换''，这里 $k^n$ 装备着对称双线性型 $(\underline{x},\underline{y})\mapsto\sum_{i=1,\ldots,n}x_iy_i$. 它与 $\mathrm{SL}_n(k)$ 的交记作 $\mathrm{SO}_n(k)$.

\emph{注}：群 $\mathrm{PSp}_{2n}(\mathbb{F}_q)$ 都是单群，而群 $\mathrm{SO}_{2n+1}(\mathbb{F}_q)$ 则是``几乎''单的.
\end{numpara}
\begin{numpara}[酉群]\label{np:fg-2-4-4}
设 $k$ 是一个配备非平凡对合 $x\mapsto\bar{x}$ 的交换域（例如：$k=\mathbb{C}$，$\overline{\cdot}$ 为复共轭；或者：$k=\mathbb{F}_{p^2}$，$\overline{\cdot}$ 为 Frobenius 映射 $x\mapsto x^p$）. 对每个整数 $n>0$，\newterm{fg-unitary-group}{酉群}{unitary group} $U_n(k)$ 是 $\mathrm{GL}_n(k)$ 中由满足 ${}^t\overline{P}.P=\mathbf{1}_n$ 的矩阵 $P$ 组成的子群. 它与 $\mathrm{SL}_n(k)$ 的交记作 $SU_n(k)$.

\emph{注}：群 $\mathrm{PSU}_n(\mathbb{F}_q)$ 都是单群.
\end{numpara}
